\paragraph{\trOrientationOperator。}\label{sub:orientation-operators}
把\cref{thm:general-surface-index-formula}\trAdapted{}到\trBrane{}的新定义上是直截了当的。
具体地说，照旧取 $(S,\nu)$，并为\trLagrange\trSubbundle{} $F_\nu\subset\partial S\times V$
配备一个\trBraneStructure。对每一条\trStriplikeEnd{} $\zeta$，由此得到一对%
\trLinear\trLagrange\trBrane{} $(\Lambda_{\zeta,0}^\#,\Lambda_{\zeta,1}^\#)$，其绝对%
\trIndex{}与\trOrientation{}空间确定 $D_{S,\nu}$ 的\trIndex{}与\trDeterminantLine，
正如\eqref{eq:general-surface-index-formula} 中那样。我们在本节的最后一项任务，是把
这一陈述扩张到稍大的一类算子上，以便能把它用于\trInhomogeneous\trPseudoholomorphicMap
的\trLinearization{}算子。
回顾\cref{sub:cauchy-riemann-operators}中定义\trNondegenerate{} $\bar\partial$-算子
所用的数据：$S$ 是一张带\trStriplikeEnd{}的\trRiemann\trSurface，$(E,\varphi_E,J_E)\to S$
是一个\trHermitian\trVectorBundle，$\nabla$ 是 $E$ 上的一个\trConnection，而
$F\subset E|\partial S$ 是一个\trLagrange\trSubbundle。它们在诸\trStriplikeEnd{}上满足
适当的\trAsymptotic{}\trConvergence{}条件。为更加实用起见，不妨给每一端添加一段区间
$\epsilon_\zeta(\{\pm\infty\}\times[0;1])$ 把诸端\trCompactification，得到一个%
\trCorneredBoundary\trSurface{} $\bar{S}$。把 $\overline{\partial S}$ 记作 $\partial S$
在 $\bar{S}$ 中的闭包，即 $\partial S$ 与诸点 $\epsilon_\zeta(\pm\infty,k)$
（$k=0,1$）的并。作为\trAsymptotic{}行为的一个后果，前所列举的诸结构在 $\bar{S}$ 上
有\trCanonical{}的扩张 $(\bar{E},\varphi_{\bar{E}},J_{\bar{E}},\bar{\nabla},\bar{F})$
（最后一项则是在 $\overline{\partial S}$ 上的扩张）；把它们限制到无穷远处的区间上时，
这些扩张恰好由我们当初称作\trLimitingDatum{}的
$(E_\zeta,\varphi_{E_\zeta},J_{E_\zeta},\nabla_\zeta,\Lambda_{\zeta,k})$ 给出。现在再
添入以下几项配料：
\begin{itemize}
\item 另一个与 $\varphi_{\bar{E}}$ \trCompatible{}的\trComplexStructure{} $I_{\bar{E}}$，
      连同\trFibrewise{}的二次复\trVolume{}形式 $\eta_{\bar{E}}^2$ 以及相伴的%
      \trSquaredPhase{}函数 $\alpha_{\bar{E}}$；
\item \trSubbundle{} $\bar{F}\subset\bar{E}|\overline{\partial S}$ 上的一个\trBraneStructure，
      由一个\trGrading{} $\bar{\alpha}^\#$ 和一个 \trPinStructure{} $\bar{P}^\#$ 组成。
\end{itemize}
任取一条\trStriplikeEnd{} $\zeta$。察看诸点 $\epsilon_\zeta(\pm\infty,k)$
（$k=0,1$），上述数据使每个 $\Lambda_{\zeta,k}$ 都成为一个\trLinear\trLagrange\trBrane{}
$\Lambda_{\zeta,k}^\#$。\trApriori{}地，它们位于互不相同的\trSymplectic\trVectorSpace
$E_{\zeta,k}$ 中，但可以用\trSymplectic\trConnection{} $\nabla_\zeta$ 的%
\trParallelTransport，连同对相关\trFiber{} $\bar{E}_z$（
$z\in\epsilon_\zeta(\{\pm\infty\}\times[0;1])$）而言的\trSquaredPhase{}映射
$\alpha_{\bar{E}}$，把 $\Lambda_{\zeta,0}^\#$ 移到 $E_{\zeta,1}$，从而定义绝对%
\trIndex{}与\trOrientation{}空间。我们断言：把这一层考虑在内，算子 $\bar\partial_\nabla$
便满足与\eqref{eq:general-surface-index-formula} 相类比的结论。
证明相当直截了当。选取一个辛\trTrivialization{} $\bar{E}\cong\bar{S}\times V$，使其在
无穷远处的区间上与诸\trConnection{} $\nabla_\zeta$ \trCompatible。其次，我们要选取
$I_{\bar{E}}$ 到常值\trComplexStructure{} $I_V$ 的一个\trDeformation，以及
$\eta_{\bar{E}}^2$ 到 $\eta_V^2$ 的一个相应的\trDeformation。对后一部分而言，
$H^1(S)$ 中有一个明显的\trObstruction，但只要原来那个\trTrivialization{}的%
\trHomotopy{}类选得恰当，它便会消失。二次复\trVolume{}形式的\trDeformation{}诱导诸%
\trSquaredPhase{}映射的一个\trHomotopy，而我们要相应地承袭诸\trGrading。结果是：
\trLagrange\trSubbundle{} $\bar{F}$——现在由一个映射
$\bar{\nu}\colon\overline{\partial S}\to Gr(V)$ 给出——作为
$S\times(V,\varphi_V,I_V)$ 的一个\trSubbundle{}带上了\trBraneStructure。最后，把 $J_{\bar{E}}$
\trDeformation{}成常值\trComplexStructure{} $I_V$，把 $\bar{\nabla}$ \trDeformation{}成平凡\trConnection。这
诱导出从 $\bar\partial_\nabla$ 到标准 \trDolbeault{}算子 $\bar\partial$ 的%
\trFredholmOperator{}空间内部的一个\trHomotopy，把我们带回到前面考虑过的特殊情形。
当 $S=H$ 时，其结果是\cref{thm:halfplane-index-determinant}的一个轻微推广：
\begin{equation}\label{eq:orientation-operator-index}
\mathrm{index}(\bar\partial_\nabla)=i(\Lambda_0^\#,\Lambda_1^\#),\qquad
\det(\bar\partial_\nabla)\cong o(\Lambda_0^\#,\Lambda_1^\#).
\end{equation}
这一情形重要到值得专有一个名字：我们称 $\bar\partial_\nabla$ 为\trLinear\trLagrange\trBrane{}对
$(\Lambda_0^\#,\Lambda_1^\#)$ 的一个\textbf{\trOrientationOperator}（orientation
operator）。
